\documentclass[10pt,a4paper]{amsart}
\usepackage[margin=19mm]{geometry}
\usepackage{amsmath,amssymb,mathtools}
\usepackage[scr=rsfs,cal=euler]{mathalfa}
\usepackage{xfrac}
\usepackage[unicode,hidelinks,bookmarks=false]{hyperref}
\newcommand{\ie}{i.e.\ }
\newcommand{\cf}{cf.\ }
\newcommand{\resp}{respectively\ }
\newcommand{\Subsection}{\subsection}
\providecommand{\nobreakdash}{\nobreak\hbox{-}}
\setlength{\emergencystretch}{2em}
\multlinegap0pt
\usepackage{amssymb}
\usepackage[shortlabels]{enumitem}
\newtheorem{theorem}{Theorem}
\title{\bf Some factorization results for formal power series}
\author{Rishu Garg, Jitender Singh}
\date{Source: arXiv:2501.05375v2 (CC BY 4.0)}
\begin{document}
\maketitle

\noindent\emph{Source and licence.} \bf Some factorization results for formal power series. Version arXiv:2501.05375v2 (CC BY 4.0), distributed by arXiv under the Creative Commons Attribution 4.0 International licence, http://creativecommons.org/licenses/by/4.0/, as recorded in the arXiv licence metadata for this exact version and shown on the abstract page https://arxiv.org/abs/2501.05375v2. Full licence notice: \texttt{licenses/arxiv-2501.05375-CC-BY-4.0.txt}.\par\smallskip

\noindent\textit{Editorial scope.} Counted unit: the article's theorem th:2G (source0.tex lines 178--181). The article's complete original proof of it is delivered with it (source0.tex lines 189--213, one \texttt{proof} environment of 25 lines, which the article places away from the statement). No mathematics is written, repaired or re-derived: the statement and the proof are the article's own lines, in the article's own notation, and no retained line is substituted or re-flowed.  No further source-local context is needed: every cross-reference of every retained line resolves inside the two delivered spans.  Nothing was omitted from any retained span, so the package carries no omission record.  No retained line contains a citation, so the wrapper carries no bibliography and no bibliography entry.  No retained line contains a figure environment, an image-inclusion command or drawing code, so no image is delivered and the package declares no asset dependency.  Editorial formatting only, in addition: an AMS \texttt{amsart} preamble that restates the article's own declarations and macros used by the retained lines, namely the environments \texttt{theorem} on separate counters, together with the standard packages those lines need.  The retained environments print in this document as Theorem 1 in order of appearance, whereas the article prints the same items with the numbering of its own full document.  The complete native source of the article as distributed in the arXiv e-print for this exact version is bundled unchanged as \texttt{sources/171402/source0.tex}.
\begin{theorem}\label{th:2G}
  Let $(R,v)$ be a discrete valuation domain with a uniformizing parameter $\pi$  for $R$. Let
  $f = \sum_{i=0}^\infty a_i z^i\in R[[z]]$ be such that $a_0= u \pi^k$ for some  unit $u\in R$ and positive integers $k$. Suppose there exists a nonnegative integer $\ell$  and an index $j\geq 1$ such that $\pi^{\ell}$ divides $a_{j}$ but  $\pi^{\ell+1}$ does not divide $a_{j}$. Then $f$ is a product of at most $\min\{k,j+\ell\}$ irreducible factors in $R[[z]]$. In particular, if $k=1$ or $j=1$ and $\ell=0$, then $f$ is irreducible.
\end{theorem}

\begin{proof}[\bf Proof of Theorem \ref{th:2G}]
Let $f(z)=f_1(z)\cdots f_r(z)$ be a product of $r$ irreducible factors $f_i$ in $R[[z]]$, $1\leq i\leq r$.  Assume without loss of generality that each $f_i$ is nonconstant. Since $f_i(0)$ is a nonzero nonunit, we have  $v(f_i(0))\geq 1$ for each $i=1,\ldots,r$. Sice $\pi$ is a uniformizing parameter for $R$, we have $v(\pi)=1$ and $f_i(0)=u_i\pi^{k_i}$ for some unit $u_i\in R$ and positive integer $k_i$ for every $i=1,\ldots,r$. Since we have
\begin{eqnarray*}
u\pi^k=f(0)=f_1(0)\cdots f_r(0),
\end{eqnarray*}
it follows that
\begin{eqnarray*}
k=v(u\pi^k)=v(f_1(0))+\cdots +v(f_r(0))=k_1+\cdots+k_r,
\end{eqnarray*}
which in view of the fact that $k_i\geq 1$ for each $i$ establishes $r\leq k$.

In the preceding paragraph, we have shown that $r\leq k$. In view of this, we may assume without loss of generality that $k\geq j+\ell$. Assume on the contrary that $r>j+\ell$, that is, $r-j> \ell$. If we write
\begin{eqnarray*}
f_i=\sum_{t=0}^{\infty}a_{it}z^t\in R[[z]],
\end{eqnarray*}
and substitute for $f_i$ in the equation $f(z)=f_1(z)\cdots f_r(z)$, then we find that
\begin{eqnarray*}
a_j &=& \sum_{i_1+i_2+\cdots+i_r=j}a_{1i_1}a_{2i_2}\cdots a_{ri_r},
\end{eqnarray*}
where the indices under the summation sign satisfy $0\leq i_1,\ldots,i_r\leq j$. Since $r-j>\ell\geq 0$, we have $r>j$. Consequently, for any given $r$-tuple $i_1,\ldots,i_r$ with $0\leq i_1,\ldots,i_r\leq j$ and $i_1+\cdots+i_r=j$, at least $r-j$ of the indices $i_1,\ldots,i_r$ is equal to 0. Since  $f_\alpha(0)=a_{\alpha0}$ and $v(f_\alpha(0))\geq 1$ for each index $\alpha\in \{1,\ldots,r\}$, it follows that  $v(a_{1i_1}a_{2i_2}\cdots a_{ri_r})\geq r-j$ for each such term. We then arrive at the following:
\begin{eqnarray*}
v(a_j)=v\Bigl(\sum_{i_1+i_2+\cdots+i_r=j}a_{1i_1}a_{2i_2}\cdots a_{ri_r}\Bigr)\geq \min_{\sum_{i_1+i_2+\cdots+i_r=j}} \Bigl\{v(a_{1i_1}a_{2i_2}\cdots a_{ri_r})\Bigr\}\geq r-j>\ell
\end{eqnarray*}
which contradicts the hypothesis.
\end{proof}

\end{document}
